\section{Can SAE Faithfulness Be Improved?}\label{sec:improve}

Section~\ref{sec:beyond} showed that the base SAEs are unfaithful at several levels. We now ask whether faithfulness
can be improved, and whether an improvement at one level carries over to the others. We tested three strategies:
fine-tuning existing SAEs on truth statements, retraining with far more data and width, and fine-tuning with a
behavioural objective. Every improved SAE was trained by us; the public \texttt{gpt2-small-res-jb} SAEs are used only
as an external check (Section~\ref{sec:e2}).

\subsection{Truth-focused fine-tuning}\label{sec:finetune}

\paragraph{A circular objective, found and corrected.} Our first truth-aware objective distilled a linear probe's
logits into the SAE and appeared to remove most of the truth gap. An audit found that the evaluation probe was the
training teacher (same rows, labels and weights), so the result was circular. In a corrected design, with the teacher
fitted on GoT training group A, evaluation probes fitted on group B and a random-teacher control, the truth-aware
objective was worse than plain reconstruction fine-tuning on every seed (TopK group-B MLP gap 0.189 versus 0.107;
random teacher 0.165). Every later fine-tuning experiment evaluates truth with probes fitted on original activations of
group B, which no training objective sees.

\paragraph{Method A across layers (IMP-L).} Method A fine-tunes an existing SAE for 500 Adam steps on GoT group-A
statements with OWT replay, with an OWT-only control that continues training for the same number of steps. After a
layer-6 pilot, the pre-registered multi-layer experiment repeated it at layers 3, 8 and 9 for all three architectures
and seeds (81 fine-tunes; \evlabel{validation}). Guards required OWT L0 within $\pm5\,\%$ of the base SAE, OWT explained
variance within 0.005 of the base, bounded splice damage and no loss of the GoT truth preference. Method A cut GoT
reconstruction error by 61--85\,\% in every cell and passed every guard. It lowered the pre-registered truth gap,
measured by a frozen nonlinear (MLP) probe fitted on original activations, on all three seeds in all nine
layer$\times$architecture cells (mean reduction 4--47\,\%), whereas the OWT-only control changed the gap by at most 0.021.
The pre-registered criterion, consistent improvement against both the base SAE and the control, was met in four of
nine cells (\evlabel{preliminary improvement}); in the other five the direction was the same but seed variation
exceeded the rule.

\paragraph{The gain is specific to one frozen reader.} Table~\ref{tab:impl_readers} compares four readers. The frozen
logistic probe, fitted on the same original activations, shows no consistent change in any cell, and neither do
re-trained logistic readers on GoT or on RepE-LM persona prompts. In the four cells where the frozen-MLP criterion was
met, the re-trained GoT gap changed by $-0.005$ to $+0.010$. Method A therefore mostly moves reconstructions back towards
the region where the original probe's decision boundary was fitted; it does not measurably restore information that a
re-trained reader can recover. This mirrors the frozen-reader fragility of Section~\ref{sec:info}. Method A also left
persona representation unchanged (every change in the frozen and re-trained persona gaps between $-0.004$ and
$+0.006$; Appendix Table~\ref{tab:rlm_finetune}).

\paragraph{Causal preservation did not robustly improve (E3).} E3 applied the M4 decomposition, with M4's code
unchanged, to the base and fine-tuned SAEs at the layer-8 final token (36 SAEs; \evlabel{validation}). The nine base
SAEs reproduce M4 exactly. Method A left causal preservation unchanged for every architecture (Method A minus base
$+0.000$, $-0.010$ and $+0.004$ for ReLU, TopK and JumpReLU, all intervals including zero), and the reconstruction still
carried only 0.61--0.69 of the effect after every fine-tune (Appendix Tables~\ref{tab:e3} and~\ref{tab:e3_classes}).

\paragraph{An output-KL term improves behaviour at a sparsity cost.} Adding a KL term between the original and
spliced output distributions to Method A lowered OWT spliced KL by 8--18\,\% and spliced $\Delta$CE by 4--19\,\% on every
seed in every cell. For ReLU and JumpReLU at layers 8 and 9, however, it also raised OWT L0 by 5.3--8.1\,\%, outside the
pre-registered guard, so those four cells are \evlabel{not creditable}; the behavioural criterion is met in the other
five. Output-KL never improved the truth gap, and its only gains in causal preservation in E3 (about $+0.02$ of the
effect) occurred in the two cells that violated the sparsity guard. TopK, whose sparsity is fixed, showed no causal
change.

% generated by scripts/shared/build_paper_artifacts.py from results/geometry_of_truth/multilayer_finetune/comparisons.csv
\begin{table}[t]
\centering
\small
\caption{IMP-L, Method A $-$ base for four readers (VALIDATION; gap change, $\downarrow$ is better; mean [95\% $t$-interval] over 3 seeds). Frozen readers are fitted once on original activations of GoT group B; re-trained readers are re-fitted on the reconstructions. $^\checkmark$: consistent (lower on all 3 seeds and $|$mean$|>2$ sd) against the base. Only the frozen MLP column is the pre-registered criterion. Source: \texttt{results/\allowbreak{}geometry\_of\_truth/\allowbreak{}multilayer\_finetune/\allowbreak{}comparisons.csv}.}
\label{tab:impl_readers}
\begin{adjustbox}{max width=\linewidth}
\begin{tabular}{llcccc}
\toprule
Layer & SAE & GoT frozen MLP & GoT frozen logistic & GoT re-trained logistic & Persona re-trained logistic \\
\midrule
3 & ReLU & -0.018 [-0.041, +0.005] & -0.000 [-0.011, +0.010] & +0.010 [-0.028, +0.048] & -0.000 [-0.007, +0.007] \\
 & TopK & -0.098 [-0.107, -0.089]$^\checkmark$ & +0.007 [-0.047, +0.062] & -0.005 [-0.104, +0.094] & +0.004 [-0.003, +0.011] \\
 & JumpReLU & -0.012 [-0.043, +0.019] & +0.005 [-0.013, +0.023] & +0.002 [-0.026, +0.031] & +0.003 [-0.001, +0.008] \\
\midrule
8 & ReLU & -0.054 [-0.095, -0.013]$^\checkmark$ & -0.003 [-0.061, +0.055] & +0.006 [-0.024, +0.035] & +0.000 [-0.002, +0.002] \\
 & TopK & -0.074 [-0.202, +0.054] & -0.045 [-0.130, +0.041] & -0.032 [-0.082, +0.019] & -0.004 [-0.012, +0.004] \\
 & JumpReLU & -0.060 [-0.247, +0.127] & -0.011 [-0.034, +0.013] & -0.005 [-0.023, +0.014] & -0.000 [-0.013, +0.012] \\
\midrule
9 & ReLU & -0.047 [-0.065, -0.030]$^\checkmark$ & +0.007 [-0.067, +0.080] & +0.010 [-0.011, +0.031] & -0.001 [-0.005, +0.003] \\
 & TopK & -0.033 [-0.061, -0.005]$^\checkmark$ & -0.007 [-0.094, +0.080] & +0.005 [-0.034, +0.043] & -0.000 [-0.003, +0.003] \\
 & JumpReLU & -0.106 [-0.162, -0.050]$^\checkmark$ & -0.002 [-0.048, +0.044] & +0.002 [-0.039, +0.043] & -0.003 [-0.014, +0.008] \\
\midrule
\multicolumn{2}{l}{consistent cells} & 5/9 & 0/9 & 0/9 & 0/9 \\
\bottomrule
\end{tabular}
\end{adjustbox}
\end{table}


\subsection{Scaling data and SAE width (R2, confirmed in M5)}\label{sec:r2}

\paragraph{Motivation and design.} The base SAEs were trained on 101{,}600 tokens at width 3{,}072, and their
validation explained variance was still rising when training stopped. It was therefore unclear which findings reflect
SAEs in general and which reflect undertraining. R2 retrained SAEs at layers 1, 8 and 11 with a modern recipe, width
12{,}288 and 100M fresh OWT tokens per SAE (about $1000\times$ more data), three seeds each, with a width ablation at
3{,}072 latents and a frozen-random-decoder baseline (Section~\ref{sec:saes}). A modern straight-through JumpReLU did not
settle at L0 = 128 in a pilot and was abandoned and disclosed before the full run. Every old and new SAE went through
the same evaluation code, which reproduces the committed results for the old SAEs, and the comparisons and their rule
were fixed and hashed before training.

\paragraph{Reconstruction and behaviour improve, and the gain is confirmed on held-out data.} All improvement-phase
comparisons were then evaluated once on the OWT and GoT test splits under a frozen protocol (M5; 81 SAEs, 552 test
rows), so these results are \evlabel{confirmed test} findings (Table~\ref{tab:m5}). Across layers and comparisons, NMSE
fell by 31--74\,\% and spliced $\Delta$CE by 35--79\,\%. For the TopK comparison at layers 8 and 11, OWT test NMSE fell
from 0.160 and 0.185 to 0.094 and 0.111, and spliced $\Delta$CE fell from 0.133 and 0.204 to 0.062 and 0.096 nats. On
validation data, equal-width retraining (more data and a better recipe alone) accounts for 23--35\,\% of the NMSE
reduction, and width for the rest. The remaining splice damage is still 1.25--1.61 times that of norm-matched random
error at layers 8 and 11 (1.45--1.96 before), so the remaining error is still more harmful than random error of the same
size \citep{engels2024darkmatter}.

\paragraph{Information preservation does not automatically follow.} No comparison improved the GoT re-trained gap or
the SAE-minus-PCA gap on test, and the width ablation made both slightly worse at layer 8 (differences of 0.012--0.015,
about a quarter of the minimum detectable effect). This is an absence of evidence within the design's sensitivity:
between-family differences smaller than about 0.05--0.09 AUROC could not have been detected, and the observed
differences are at most 0.025. On validation data, the re-trained persona gap fell by 27--31\,\% at layers 8 and 11 for
TopK, but the persona gap relative to error-matched PCA did not improve, and at layer 8 the random-decoder baseline had a
smaller GoT re-trained gap than the trained SAE (0.052 versus 0.081), so that metric does not discriminate there.

\paragraph{Causal preservation improves on validation only.} For TopK, the share of the persona effect carried by the
reconstruction rose from 0.62 to 0.73 at layer 8 and from 0.62 to 0.69 at layer 11 (\evlabel{validation}). The layer-8
gain passes the pre-registered seed-level rule but is close to the per-SAE minimum detectable effect at that site, and
RepE-LM had no evaluated test split, so it is not a confirmed result. Scaling therefore removes a large part of the
behavioural unfaithfulness we measured, and probably part of the causal shortfall, but it does not by itself make the
SAE preserve truth information better.

% generated by scripts/shared/build_paper_artifacts.py from results/shared/m5_confirmatory_test/comparisons.csv
\begin{table}[t]
\centering
\small
\caption{M5: the single pre-registered confirmatory TEST pass (OWT test, 100 sequences, 12,700 tokens; GoT cities+neg\_cities test). Same rules as on validation (C/A: all 9 seed pairs and Welch CI; K: all 3 paired seeds and $|\bar d|>2\,$sd). Means over 3 seeds. C1: TopK old-3072 $\rightarrow$ R2-12288; C3: old JumpReLU-3072 $\rightarrow$ R2 BatchTopK-12288; K1: R4 KL fine-tune vs MSE control. All comparisons are counted in Table~\ref{tab:m5_counts}. Source: \texttt{results/\allowbreak{}shared/\allowbreak{}m5\_confirmatory\_test/\allowbreak{}comparisons.csv}.}
\label{tab:m5}
\begin{adjustbox}{max width=\linewidth}
\begin{tabular}{lllcccl}
\toprule
ID & Layer & Metric & Before & After & Improvement & Class (TEST) \\
\midrule
C1 & L1 & OWT NMSE $\downarrow$ & 0.0787 & 0.0248 & +68.6\% & IMPROVEMENT \\
 & L1 & spliced $\Delta$CE $\downarrow$ & 0.0221 & 0.0056 & +74.9\% & IMPROVEMENT \\
 & L8 & OWT NMSE $\downarrow$ & 0.1598 & 0.0941 & +41.1\% & IMPROVEMENT \\
 & L8 & spliced $\Delta$CE $\downarrow$ & 0.1330 & 0.0620 & +53.4\% & IMPROVEMENT \\
 & L8 & GoT re-fit gap $\downarrow$ & 0.0607 & 0.0558 & +8.1\% & no clear change \\
 & L11 & OWT NMSE $\downarrow$ & 0.1848 & 0.1112 & +39.8\% & IMPROVEMENT \\
 & L11 & spliced $\Delta$CE $\downarrow$ & 0.2041 & 0.0962 & +52.9\% & IMPROVEMENT \\
 & L11 & GoT re-fit gap $\downarrow$ & 0.1009 & 0.0775 & +23.2\% & no clear change \\
\midrule
C3 & L1 & OWT NMSE $\downarrow$ & 0.0675 & 0.0178 & +73.7\% & IMPROVEMENT \\
 & L1 & spliced $\Delta$CE $\downarrow$ & 0.0088 & 0.0019 & +78.8\% & IMPROVEMENT \\
 & L8 & OWT NMSE $\downarrow$ & 0.1560 & 0.0897 & +42.5\% & IMPROVEMENT \\
 & L8 & spliced $\Delta$CE $\downarrow$ & 0.1212 & 0.0542 & +55.3\% & IMPROVEMENT \\
 & L8 & GoT re-fit gap $\downarrow$ & 0.0584 & 0.0654 & -12.0\% & no clear change \\
 & L11 & OWT NMSE $\downarrow$ & 0.1851 & 0.1073 & +42.0\% & IMPROVEMENT \\
 & L11 & spliced $\Delta$CE $\downarrow$ & 0.2155 & 0.0927 & +57.0\% & IMPROVEMENT \\
 & L11 & GoT re-fit gap $\downarrow$ & 0.1070 & 0.0871 & +18.6\% & no clear change \\
\midrule
K1 & L8 & OWT NMSE $\downarrow$ & 0.0938 & 0.0958 & -2.1\% & \textbf{REGRESSION} \\
 & L8 & spliced $\Delta$CE $\downarrow$ & 0.0623 & 0.0514 & +17.4\% & IMPROVEMENT \\
 & L8 & GoT re-fit gap $\downarrow$ & 0.0585 & 0.0484 & +17.3\% & no clear change \\
 & L11 & OWT NMSE $\downarrow$ & 0.1109 & 0.1133 & -2.2\% & \textbf{REGRESSION} \\
 & L11 & spliced $\Delta$CE $\downarrow$ & 0.0964 & 0.0807 & +16.3\% & IMPROVEMENT \\
 & L11 & GoT re-fit gap $\downarrow$ & 0.0780 & 0.0784 & -0.5\% & no clear change \\
\bottomrule
\end{tabular}
\end{adjustbox}
\end{table}


\subsection{Behaviour-targeted training at fixed sparsity (R4)}\label{sec:r4}

The output-KL results above were confounded by sparsity drift. R4 therefore fine-tuned the strong R2 TopK-12288 SAEs,
whose L0 is fixed at exactly 128, with an NMSE plus output-KL objective \citep{karvonen2025kl,braun2024e2e} for 5M
tokens (about 5\,\% of the training budget) at layers 8 and 11, against an MSE-only control on identical tokens. On the
held-out test (M5, \evlabel{confirmed test}), spliced $\Delta$CE fell from 0.062 and 0.096 to 0.051 and 0.081 nats at
layers 8 and 11, an improvement of 16--17\,\%. The cost is a regression in reconstruction: test NMSE rose by about 2\,\%
(0.094 to 0.096 and 0.111 to 0.113), within the pre-registered explained-variance guard but classified as a regression
in NMSE. On validation data (Appendix Table~\ref{tab:r4}), no information gap changed, the per-pair persona
direction cosine dropped by 0.012--0.013, and causal preservation at layer 11 \emph{decreased} from 0.694 to 0.671 on
every seed (\evlabel{regression}); at layer 8 it did not change. A behavioural objective therefore improves behaviour
creditably once sparsity is fixed, but it does so as a trade-off: a small reconstruction cost, no information gain and
a loss of causal preservation at one of two sites.

\subsection{External validation with a public SAE (E2)}\label{sec:e2}

A natural concern is that our conclusions describe small, in-house SAEs. E2 applied the same evaluation to the public
\texttt{gpt2-small-res-jb} SAEs (24{,}576 latents, 300M tokens) at the pre-registered layers 8 and 10 and the
extension layers 1 and 11 (\evlabel{validation}; Appendix Table~\ref{tab:e2}). The broad pattern holds. The frozen GoT
loss is mostly reader fragility; the persona loss is real and survives re-training; the SAE does not beat PCA at equal
reconstruction error; and the reconstruction carries only part of the persona effect at layers 8 and 11 (0.56 and
0.61). The one exception is the GoT re-trained gap at layer 8, which is small and has an interval touching zero, as for
our own SAEs. This is one public family with a ReLU architecture, so it shows that the conclusions are not an artefact
of our training recipe, not that they hold for every SAE.
